\documentclass[preview]{standalone}

\usepackage{amsmath}
\usepackage{amssymb}
\usepackage{bettelini}
\usepackage{stellar}
\usepackage{definitions}

\begin{document}

\id{measuretheory-vitali-sets}
\genpage

\section{Vitali Sets}

\begin{snippetexample}{vitali-non-lebesgue-measurable-set}{A non-Lebesgue-measurable set}
    Assume the axiom of choice. On \([0,1]\), define
    \[
        x\sim y
        \quad\Longleftrightarrow\quad
        x-y\in\rationalnumbers.
    \]
    Each equivalence class is countable. Choose one representative from every
    class and let \(V\) be the resulting set. It has the cardinality of the
    continuum and contains exactly one rational number.

    For \(r\in\rationalnumbers\cap[0,1]\), translate \(V\) modulo one:
    \[
        V_r=
        \bigl((V+r)\cap[0,1]\bigr)
        \cup
        \bigl((V+r-1)\cap[0,1]\bigr).
    \]
    The two pieces are disjoint. The sets \(V_r\) are pairwise disjoint and
    cover \([0,1]\). Indeed, every \(x\in[0,1]\) is equivalent to exactly one
    representative \(v\in V\), so \(x-v\) is rational modulo one. If two
    translates met, two representatives of \(V\) would differ by a rational,
    contradicting their choice from distinct classes.

    If \(V\) were Lebesgue measurable, every \(V_r\) would be measurable and
    translation invariance would give \(\mu(V_r)=\mu(V)\). If
    \(\mu(V)=0\), countable additivity would imply
    \[
        1=\mu([0,1])
        =\sum_{r\in\rationalnumbers\cap[0,1]}\mu(V_r)=0.
    \]
    If \(\mu(V)>0\), the same sum would be \(+\infty\). Both alternatives are
    impossible. Hence \(V\) is not Lebesgue measurable.
\end{snippetexample}

\begin{snippetproposition}{lebesgue-outer-measure-not-finitely-additive}{Lebesgue outer measure is not finitely additive}
    There are finitely many pairwise disjoint subsets
    \(A_1,\ldots,A_N\subseteq[0,1]\) such that
    \[
        \mu^*\!\left(\bigcup_{n=1}^NA_n\right)
        <
        \sum_{n=1}^N\mu^*(A_n).
    \]
\end{snippetproposition}

\begin{snippetproof}{lebesgue-outer-measure-not-finitely-additive-proof}{lebesgue-outer-measure-not-finitely-additive}{Lebesgue outer measure is not finitely additive}
    Use the modulo-one translates \(V_r\) from the Vitali construction.
    Piecewise translation modulo one preserves outer measure, so every
    \(V_r\) has the same outer measure as \(V\). Subadditivity and
    \(\bigcup_rV_r=[0,1]\) give
    \[
        1
        =\mu^*([0,1])
        \leq
        \sum_{r\in\rationalnumbers\cap[0,1]}\mu^*(V_r).
    \]
    Thus their common outer measure is positive. Choose distinct
    \(r_1,\ldots,r_N\) so that
    \[
        \sum_{n=1}^N\mu^*(V_{r_n})>1.
    \]
    Since their union is contained in \([0,1]\),
    \[
        \mu^*\!\left(\bigcup_{n=1}^NV_{r_n}\right)
        \leq1
        <
        \sum_{n=1}^N\mu^*(V_{r_n}).
    \]
    The missing equality is precisely the failure of finite additivity.
\end{snippetproof}

\begin{snippetexercise}{nonmeasurable-set-in-euclidean-space-exercise}{A nonmeasurable set in higher-dimensional Euclidean space}
    Starting from a Vitali set \(V\subseteq[0,1]\), construct a
    non-Lebesgue-measurable subset of \(\realnumbers^n\). One may use
    \(V\cartesianprod[0,1]^{n-1}\) and the section form of the completed
    Fubini theorem. Notice that the Cartesian product of two nonmeasurable
    sets need not automatically be nonmeasurable.
\end{snippetexercise}

\begin{snippetexercise}{positive-measure-set-contains-nonmeasurable-subset-exercise}{Nonmeasurable subsets of positive-measure sets}
    Prove that every Lebesgue measurable subset of \(\realnumbers\) with
    positive measure contains a non-Lebesgue-measurable subset. The set need
    not contain any interval.
\end{snippetexercise}

\end{document}
